\emph{Conventions and hypotheses.} As in Theorem \claimref{synthesis.field-strength-cancels-from-the-leading-null-singularity-of-hadamard-dirac-states-on-null-hypersurfaces}: physical vectors as in Fact \claimref{synthesis.imported-gtg-dictionary}, signature $(+,-,-,-)$; $\mathcal{H}$ is generated by the affinely parametrised null geodesics $u \mapsto X(u, y)$ with tangent $k$; the null frame $\{k, n, f_{1}, f_{2}\}$ is parallel along each generator, $k\cdot n = 1$; $J$, $a = |\det J|$, $B = J'J^{-1}$ (self-adjoint, $\mathcal{H}$ being twist-free), $\theta = \mathrm{tr}\,B = a'/a$, the Jacobi map $A$, $M(u, u') = (u - u')\bar J(u')A^{-1}J(u)$ and $\Delta = (u - u')^{2}/\det A$ on the generator are those of that theorem's (i). The piece of $\mathcal{H}$ considered lies in a geodesically convex set $N$ and satisfies (S): points of distinct generators are spacelike separated, $\sigma(X(u, y), X(u', y')) < 0$ for $y \neq y'$. (S) holds on every small enough piece: by that theorem's (i), $\sigma(X(u, y' + \eta), X(u', y')) = -\tfrac12 M(u, u')(\eta, \eta) + O(|\eta|^{3})$ uniformly for $(u, u')$ in a compact interval, and $M(u, u') \to \bar J J$ as $u' \to u$ is positive, so $\sigma < 0$ for $0 < |\eta| < \delta_{0}$ and $|u - u'| < \delta_{0}$, and a piece of transverse and generator extent below $\delta_{0}$ has (S). (S) also holds on every local Rindler horizon of Premise \claimref{synthesis.c-loc-local-horizon-premise}: the null boundary of the past of a small two-surface element is achronal, and in a small caustic-free patch two null-related points of it lie on one generator (standard causal theory, cited from memory). The parametrix is $G_{\mathrm{par}} = (iD_{x} + m_{0})H_{+}e_{0}$, $H = U/\sigma + V\log\sigma$ with $U = \Delta^{1/2}R$ and $V$ smooth (Fact \claimref{synthesis.imported-hadamard-parametrix-of-the-squared-dirac-operator-leads-with-van-vleck-root-times-spin-transport}(a), (b)), in the form of Fact \claimref{synthesis.imported-dirac-hadamard-form-propagates-and-adiabatic-states-are-hadamard}(a) (the Dirac operator with its mass term applied to a Hadamard parametrix of its square); its boundary value is taken with $\sigma_{\epsilon}(x, x') = \sigma(x, x') + i\epsilon(t(x) - t(x')) - \tfrac12\epsilon^{2}$ for a time function $t$ (Fact \claimref{synthesis.imported-hadamard-parametrix-of-the-squared-dirac-operator-leads-with-van-vleck-root-times-spin-transport}(c); the sign of $i$ is the one that gives the $+i0$ of $W_{0}$ in the flat computation of Theorem \claimref{synthesis.bulk-hadamard-states-restrict-to-the-generator-condition-on-flat-null-planes}(ii)). Only $\mathrm{Re}\,\sigma_{\epsilon} = \sigma - \tfrac12\epsilon^{2} \le \sigma$ and $\mathrm{Im}\,\sigma_{\epsilon} = \epsilon s$, $s = t(x) - t(x')$, are used below.

\emph{Step 0 (adapted rotation gauge).} A rotation gauge transformation (Definition \claimref{obj.synthesis.gtg-flat-spacetime-gauge-fields}) makes the images of the parallel frame the constant vectors $k = e_{0} + e_{3}$, $n = \tfrac12(e_{0} - e_{3})$, $f_{a} = e_{a}$ at every point of the piece of $\mathcal{H}$ (any smooth rotor field on $\mathcal{H}$ extends smoothly off it). Then $\Lambda = \tfrac12 nk$ is constant on $\mathcal{H}$ and $\Lambda^{\dagger} = \Lambda$, and since $D_{V}e = \Omega(V)\cdot e$ for a frame vector with constant image: $\Omega(k) = 0$ on $\mathcal{H}$ (the frame is parallel along generators), and for $V$ tangent to $\mathcal{H}$, $\Omega(V)\cdot k = D_{V}k$, whose screen part is $B(V_{\perp})$ ($D_{E}k = D_{k}E$ for the deviation fields $E = \partial_{y}X$). So the $f\wedge n$ part of $\Omega(V)$ is $B(V_{\perp})\wedge n$. In another rotation gauge spinors change by a rotor $S(x)$, the two-point kernel by $S(x)\,(\cdot)\,S(x')^{\dagger}$, and the projector by $S\Lambda S^{-1}$ on the left and $(S\Lambda S^{-1})^{\dagger}$ on the right; the components of the restricted kernel in the parallel frame do not change (computed exactly in the model below, where the gauge is not adapted off one generator and the right projector is $\Lambda(x')^{\dagger} = e_{0}\tilde\Lambda(x')e_{0}$).

\emph{Step 1 (decomposition).} Write $x = X(u, y' + \eta)$, $x' = X(u', y')$. For scalar $\phi$ and spinor-valued $U$, $D(\phi U) = (\nabla\phi)U + \phi\,DU$, so with $\nabla\sigma_{\epsilon} = \nabla\sigma + i\epsilon\nabla t$,
\[
\Lambda D\Bigl(\frac{U}{\sigma_{\epsilon}}\Bigr)e_{0}\Lambda = \frac{N}{\sigma_{\epsilon}^{2}} + \frac{(i\epsilon s - \tfrac12\epsilon^{2})\,\Lambda(DU)e_{0}\Lambda}{\sigma_{\epsilon}^{2}} - \frac{i\epsilon\,\Lambda(\nabla t)Ue_{0}\Lambda}{\sigma_{\epsilon}^{2}},\qquad N = \Lambda\bigl[(DU)\sigma - (\nabla\sigma)U\bigr]e_{0}\Lambda,
\]
and the mass and logarithmic parts are $m_{0}\Lambda Ue_{0}\Lambda/\sigma_{\epsilon}$ and $\Lambda(iD + m_{0})(V\log\sigma_{\epsilon})e_{0}\Lambda$. In the null frame $D = nD_{k} + kD_{n} + f^{c}D_{f_{c}}$ ($f^{c} = -f_{c}$), and $\Lambda k = 0$ removes $kD_{n}$ entirely. Everything is smooth in $(u, u', y', \eta)$ on $N\times N$ except through $1/\sigma_{\epsilon}$ and $\log\sigma_{\epsilon}$.

\emph{Step 2 (no terms of degree $-4$ or $-3$; proved).} On the generator ($\eta = 0$) $\sigma = 0$ and $\nabla\sigma = (u - u')k$ for all $u, u'$, and $\sigma$ has no term linear in $\eta$ (that theorem's (i)); hence $\partial_{u}\sigma = k\cdot\nabla\sigma = O(|\eta|^{2})$. Expand $N = N_{0} + N_{1} + N_{2} + O(|\eta|^{3})$ in homogeneous degrees in $\eta$, $U = \Delta^{1/2}(1 + R_{1} + \ldots)$ with $R_{1}$ a bivector linear in $\eta$ ($R = 1$ on the generator). $N_{0} = -(u - u')\Delta^{1/2}\Lambda ke_{0}\Lambda = 0$. Since $(DU)\sigma = O(|\eta|^{2})$, $N_{1} = -\Lambda\bigl[w\Delta^{1/2} + (u - u')k(\ldots)\bigr]e_{0}\Lambda$, $w$ the linear part of $\nabla\sigma$; the second term vanishes by $\Lambda k = 0$ and the first is $-\Delta^{1/2}(w\cdot k)\Lambda$ by $\Lambda pe_{0}\Lambda = (p\cdot k)\Lambda$ (Theorem \claimref{synthesis.bulk-hadamard-states-restrict-to-the-generator-condition-on-flat-null-planes}(ii)), and $w\cdot k$ is the linear part of $\partial_{u}\sigma$, zero. So $N = O(|\eta|^{2})$ and $N/\sigma_{\epsilon}^{2}$ has no part of degree below $-2$. Likewise $\Lambda Ue_{0}\Lambda = \Delta^{1/2}\Lambda e_{0}\Lambda + O(|\eta|) = O(|\eta|)$ ($\Lambda e_{0}\Lambda = 0$) and $\Lambda\nabla\sigma = (u - u')\Lambda k + O(|\eta|) = O(|\eta|)$, so the mass term $m_{0}\Lambda Ue_{0}\Lambda/\sigma$ and the term $\Lambda(\nabla\sigma)Ve_{0}\Lambda/\sigma$ of $D(V\log\sigma)$ have degree $\ge -1$; $(\partial_{u}\sigma)V/\sigma$ has degree $0$ and $(DV)\log\sigma$ is integrable.

\emph{Step 3 (the coefficient of $\log\epsilon$; proved).} For a function $g$ of $\eta$ homogeneous of degree $-2$ let $L[g] = \lim_{\delta \to 0}\bigl(\int_{\delta < |\eta|_{M} < 1}g\,d^{2}\eta\bigr)/\log(1/\delta)$, $|\eta|_{M}^{2} = M(\eta, \eta)$. With $\zeta = L_{M}\eta$, $\bar L_{M}L_{M} = M$, $L[g] = (\det M)^{-1/2}\oint g(L_{M}^{-1}\omega)\,d\omega$; for $g = \eta^{T}X\eta/\sigma_{2}^{2}$, $\sigma_{2} = -\tfrac12 M(\eta, \eta)$, this is $4\pi\,\mathrm{tr}(M^{-1}X)/\sqrt{\det M}$, and $L[1/\sigma_{2}] = -4\pi/\sqrt{\det M}$. The regularised integral has the cutoff at $|\eta|_{M}^{2} \sim \epsilon|s|$, so $\int F(y' + \eta)\,g_{\epsilon}\,d^{2}\eta = -\tfrac12 L[g]F(y')(\log\epsilon + \log(-is)) + O(1)$ (for $g = 1/\sigma_{2}$ this is the transverse-integral formula of that theorem's Step 5). The degree $-2$ part of $N/\sigma^{2}$ is $N_{2}/\sigma_{2}^{2}$, with three contributions.
(a) From $nD_{k}$: $\Lambda nXe_{0}\Lambda = X\Lambda$ for scalar $X$ ($\Lambda n = n$, $ne_{0}\Lambda = \Lambda$), and the $R_{1}$ parts are of degree $\ge -1$; so the contribution is the degree $-2$ part of $\partial_{u}(\Delta^{1/2}/\sigma)\Lambda$, and since $\partial_{u}$ acts on the coefficients, $L$ of it is $\partial_{u}L[\Delta^{1/2}/\sigma_{2}] = -4\pi\,\partial_{u}(\Delta^{1/2}/\sqrt{\det M})$.
(b) From the connection in $f^{c}D_{f_{c}}$: $\tfrac12\Lambda f^{c}\Omega(f_{c})Ue_{0}\Lambda/\sigma = \tfrac12\theta(u)\Delta^{1/2}\Lambda/\sigma + O(|\eta|^{-1})$, since only the $f\wedge n$ part of a bivector survives $\Lambda f^{c}(\cdot)e_{0}\Lambda$, with $\Lambda f^{c}(w\wedge n)e_{0}\Lambda = f^{c}w\Lambda$, and $f^{c}B(f_{c}) = \mathrm{tr}\,B$ for self-adjoint $B$ (exact Clifford identities; the shear and the other components drop). $L$ of it is $-2\pi\theta\,\Delta^{1/2}/\sqrt{\det M}$.
(c) From the derivative in $f^{c}D_{f_{c}}$: $\Lambda f^{c}\partial_{f_{c}}(U/\sigma)e_{0}\Lambda$. Its scalar part $\Lambda f^{c}\partial_{f_{c}}(\Delta^{1/2}/\sigma)e_{0}\Lambda$ vanishes ($\Lambda f^{c}e_{0}\Lambda = 0$), and what remains at degree $-2$ is $\Lambda f^{c}\partial_{f_{c}}(\Delta^{1/2}R_{1}/\sigma)e_{0}\Lambda$, a transverse derivative of a function homogeneous of degree $-1$. With $E_{b} = J_{cb}f_{c} + \alpha_{b}k$, $\partial_{f_{c}} = (J^{-1})_{bc}(\partial_{\eta_{b}} - \alpha_{b}\partial_{u})$; the $\alpha\partial_{u}$ part and the $\eta$-dependence of $J$ give degree $\ge -1$, so at degree $-2$ this is the divergence $\partial_{\eta_{b}}V_{b}$ of $V_{b} = (J^{-1})_{bc}\Delta^{1/2}\Lambda f^{c}R_{1}e_{0}\Lambda/\sigma_{2}$, homogeneous of degree $-1$. Its integral over $\delta < |\eta|_{M} < 1$ is a difference of two boundary integrals, each independent of $\delta$ by homogeneity: $L = 0$. The first-order transport therefore does not enter, whatever $R_{1}$ is.
Hence the coefficient of $\log\epsilon$ (and of $\log(-is)$) in the restriction of the leading term is, with the canonical factor,
\[
-\tfrac12 L[N_{2}/\sigma_{2}^{2}]\,(a(u)a(u'))^{1/2} = 2\pi\,(a(u)a(u'))^{1/2}\Bigl(\partial_{u} + \tfrac12\theta(u)\Bigr)\frac{\Delta^{1/2}}{\sqrt{\det M}}\,\Lambda\;\cdot F(y'),
\]
times the normalising constant (computed exactly on generic jets, below). By Theorem \claimref{synthesis.field-strength-cancels-from-the-leading-null-singularity-of-hadamard-dirac-states-on-null-hypersurfaces}(i), $\Delta^{1/2}/\sqrt{\det M} = (a(u)a(u'))^{-1/2}$, and $(\partial_{u} + \tfrac12\theta(u))(a(u)a(u'))^{-1/2} = 0$ because $\theta = a'/a$. The coefficient vanishes identically in $(u, u', y')$: the term $-\pi\theta(u)(\log\epsilon + \log(-is))\delta^{2}(y - y')$ of that theorem's Step 5, from (a), is cancelled by the projected transverse connection (b), and there is no $\log(u - u')$ term in the restricted kernel.

\emph{Step 4 (model integrals; proved).} Near the generator write $\sigma_{\epsilon} = -A_{\epsilon} + \rho$ with the model $A_{\epsilon}(\eta) = \tfrac12 M(\eta, \eta) + c_{0}\epsilon^{2} - i\epsilon(s_{0} + s_{1}\cdot\eta)$, where $s_{0} = t(X(u, y')) - t(X(u', y'))$, $s_{1}$ is the linear part of $s$ in $\eta$, $c_{0} = \tfrac12$, and $|\rho| \le C(|\eta|^{3} + \epsilon|\eta|^{2} + \epsilon^{2}|\eta|)$ (the $\epsilon^{2}$ coefficient of $\sigma_{\epsilon}$ is $\tfrac12$ here; for another standard regularisation it is a smooth positive $c(x, x')$, and $c_{0} = c$ on the generator). Since $\mathrm{Re}\,A_{\epsilon} \ge \tfrac12 M(\eta, \eta) + c_{0}\epsilon^{2}$, also along the shift $\eta \to \eta + i\lambda\epsilon M^{-1}s_{1}$, $\lambda \in [0, 1]$ (the real part gains $\epsilon^{2}(\lambda - \tfrac12\lambda^{2})s_{1}\cdot M^{-1}s_{1} \ge 0$), the shift of the integration plane crosses no zero, and at $\lambda = 1$ the model becomes isotropic in $\zeta$: $A_{\epsilon} = \tfrac12|\zeta|^{2} + C_{\epsilon}$, $C_{\epsilon} = g\epsilon^{2} - i\epsilon s_{0}$, $g = c_{0} + \tfrac12 s_{1}\cdot M^{-1}s_{1} > 0$. With $\int_{0}^{\infty}dq/(q + C)^{2} = 1/C$ this gives, exactly:
(a) $\int_{\R^{2}}d^{2}\eta/A_{\epsilon}^{2} = (2\pi/\sqrt{\det M})/C_{\epsilon}$;
(b) $\int_{\R^{2}}(i\epsilon(s_{0} + s_{1}\cdot\eta) - c_{0}\epsilon^{2})/A_{\epsilon}^{2}\,d^{2}\eta = -(2\pi/\sqrt{\det M})(1 + \tfrac12\epsilon^{2}s_{1}\cdot M^{-1}s_{1}/C_{\epsilon})$, bounded by $(2\pi/\sqrt{\det M})(1 + s_{1}\cdot M^{-1}s_{1}/(2g))$ (the numerator is $\tfrac12 M(\eta, \eta) - A_{\epsilon}$; the odd term of the shift integrates to zero);
(c) for $X$ with $\mathrm{tr}(M^{-1}X) = 0$ and $\hat A_{\epsilon} = \tfrac12 M(\eta, \eta) + c_{0}\epsilon^{2} - i\epsilon s_{0}$: $\int_{|\eta|_{M} < \rho}\eta^{T}X\eta/\hat A_{\epsilon}^{2} = 0$ and $\int_{|\eta|_{M} < \rho}\eta^{T}X\eta\,(s_{1}\cdot\eta)/\hat A_{\epsilon}^{3} = 0$ (angular means), while $|\eta^{T}X\eta|\,|A_{\epsilon}^{-2} - \hat A_{\epsilon}^{-2} - 2i\epsilon(s_{1}\cdot\eta)\hat A_{\epsilon}^{-3}| \le C\epsilon^{2}|\eta|^{4}(|\eta|^{2} + \epsilon^{2})^{-4}$, whose integral is bounded uniformly in $\epsilon$ and $s_{0}$.

\emph{Step 5 (the $i0$ term and the bounded remainder; proved).} Fix profiles $F$, $G$ with support in a compact set $\mathcal{K}$ of labels and $\|F\|_{C^{1}}, \|G\|_{C^{1}} \le 1$, and $(u, u')$ in a compact interval. Split the $\eta$ integration at $|\eta| = \delta_{0}$. For $|\eta| \ge \delta_{0}$, (S) and compactness give $\sigma \le -c_{1} < 0$, so $|\sigma_{\epsilon}| \ge c_{1}$ and the integrand is smooth, bounded uniformly in $\epsilon$ and uniformly convergent. For $|\eta| < \delta_{0}$, $|\sigma_{\epsilon}|, |A_{\epsilon}| \ge c(|\eta|^{2} + \epsilon^{2})$ uniformly; each term of Step 1 is its model (coefficients frozen at $\eta = 0$, denominators replaced by $A_{\epsilon}$) plus a correction bounded by $C|\eta|^{-1}$ plus finitely many terms $C\epsilon^{i}|\eta|^{j}(|\eta|^{2} + \epsilon^{2})^{-(i + j + 2)/2}$ with $i \ge 1$, $j \ge 0$, all with integrals over $|\eta| < \delta_{0}$ bounded uniformly in $\epsilon$ (scale $\eta = \epsilon\xi$; the condition $i \ge 1$ makes the $\xi$ integral converge), using $F(y' + \eta) - F(y') = O(|\eta|)$, $N - N_{2} = O(|\eta|^{3})$, the bound on $\rho$ and the $O(|\eta|)$ variation of $\Delta^{1/2}$, $\Lambda(\nabla t)e_{0}\Lambda$ and $\Lambda(DU)e_{0}\Lambda$. The models:
(1) $N/\sigma_{\epsilon}^{2}$: model $F(y')N_{2}/A_{\epsilon}^{2}$, bounded uniformly by Step 4(c), since $\mathrm{tr}(M^{-1}X) = 0$ for the quadratic form $X$ of $N_{2}$ (Step 3).
(2) The middle term: model $F(y')\Lambda(DU)e_{0}\Lambda|_{\eta = 0}$ times the integral of Step 4(b): bounded.
(3) The $i0$ term: at $\eta = 0$, $\Lambda(\nabla t)Ue_{0}\Lambda = \Delta^{1/2}(k\cdot\nabla t)\Lambda = \Delta^{1/2}\partial_{u}t\,\Lambda$, and by Step 4(a) the model is $(2\pi\Delta^{1/2}/\sqrt{\det M})\,\partial_{u}t/(s_{0} + ig\epsilon)\,\Lambda F(y')$ (the region $|\eta| \ge \delta_{0}$ of the model is $O(\epsilon)$). Since $t$ increases strictly along the future-directed generators, $s_{0} = (u - u')\varphi(u, u')$ with $\varphi > 0$ smooth, and as $\epsilon \to 0$, $\partial_{u}t/(s_{0} + ig\epsilon) \to \partial_{u}\log(s_{0} + i0) = (u - u' + i0)^{-1} + \partial_{u}\log\varphi$ in $\mathcal{D}'$ ($g > 0$ smooth). With the canonical factor, $(a(u)a(u'))^{1/2}\Delta^{1/2}/\sqrt{\det M} = 1$, and the constant fixed by the flat case, this is $W_{0}(u, u')\Lambda$ smeared with $\langle F, G\rangle$ plus the smooth kernel $\tfrac{i}{2\pi}\partial_{u}\log\varphi\,\Lambda$, which depends on the time function (on the flat null plane with $t$ nonlinear along the generators the limit is $\tfrac{i}{2\pi}\partial_{u}t/(t(u) - t(u') + i0)$, not $W_{0}$; the bulk time-function independence of the Fact does not pass to the restriction), so the regularised restriction is defined for the chosen $t$ and only modulo smooth kernels: that theorem's (ii), with the exclusion of every other term of this order now proved (Steps 2, 3: nothing of degree below $-2$, no $\log$ at degree $-2$, and the only source of $(u - u')^{-1}$ is the $i0$ term).
(4) The mass and logarithmic terms: degree $\ge -1$ by Step 2, so bounded by $C|\eta|^{-1}$ or $C(1 + |\log|\sigma_{\epsilon}||)$, integrable uniformly; the term $i\epsilon\Lambda(\nabla t)Ve_{0}\Lambda/\sigma_{\epsilon}$ is $O(\epsilon\log(1/\epsilon))$.
For fixed $u \neq u'$ every correction and every model other than (3) converges as $\epsilon \to 0$ (where needed, $|\sigma_{\epsilon}| \ge \tfrac12\epsilon|s_{0}|$ near $\eta = 0$ improves the bounds to $O(\epsilon^{1/2}|s_{0}|^{-1/2})$). By dominated convergence in $(u, u')$, the smeared restricted parametrix converges in $\mathcal{D}'$ to $W_{0}\langle F, G\rangle\Lambda + r_{F,G}$ with $\sup|r_{F,G}| \le C$, $C$ depending on $\mathcal{K}$, the interval, $m_{0}$ and the $C^{1}$ bounds only. This is (ii) of the claim.

\emph{Step 6 (the bound of (HS); proved).} For $r$ bounded, $|\lambda^{\varpi}\int r(u_{0} + \lambda x, u_{0} + \lambda x')\phi(x, x')\,dx\,dx'| \le \lambda^{\varpi}\sup|r|\,\|\phi\|_{L^{1}} \to 0$ for every $\varpi > 0$, uniformly over the profiles of Step 5. So $R^{\mathrm{par}}$ has generator scaling degree at most $0$ at every $(u_{0}, u_{0})$, and the clause of Premise \claimref{synthesis.hadamard-scaling-limit-premise} holds with any $\varpi \in (0, 1)$, for every $m_{0} \ge 0$ and every geodesic null hypersurface, in particular every local Rindler horizon through $x_{0}$ (a smooth null hypersurface with affinely parametrised generators, Premise \claimref{synthesis.c-loc-local-horizon-premise}), whatever the field strength at $x_{0}$. Nothing in Steps 0 to 6 uses vanishing expansion or shear, the field equations, or the value of the field strength.

\emph{Exact checks (computed).} (1) The projector identities of Steps 2 and 3 in the spacetime algebra. (2) Generic jets in the adapted gauge: $M(u)$, $\partial_{u}M$, cubic terms of $\sigma$, $\Delta^{1/2}$ with transverse dependence, a transport rotor with arbitrary first- and second-order parts (all bivector components), a connection with the $f\wedge n$ part $B(f_{c})\wedge n$ plus arbitrary other components and transverse dependence, $J$ depending on $\eta$ and arbitrary $\alpha$: $N_{0} = N_{1} = 0$ and $\mathrm{tr}(M^{-1}X) = -[\partial_{u}d - \tfrac12 d\,\mathrm{tr}(M^{-1}\partial_{u}M) + \tfrac12\theta d]\Lambda$ ($d = \Delta^{1/2}$ on the generator) exactly, at four sample sets; with $B$ not self-adjoint or the connection dropped the formula fails. (3) On the plane-wave model of that theorem (tidal map rotated off the label basis, deviation map not commuting with it, Ricci and Weyl parts nonzero), $\partial_{u}\log\Delta - \mathrm{tr}(M^{-1}\partial_{u}M) + \theta(u) = 0$ exactly. (4) From first principles on that field in its Kerr-Schild gauge $\hb^{-1}(a) = a + \tfrac12 H(a\cdot K)K$ (not adapted off the central generator): with the closed-form Synge function, the exact transport rotor $R = 1 + \tfrac12\sum_{i}(t_{i} - w_{i})e_{i}\wedge K$ (the null rotations commute; $t$, $w$ the end velocities of the geodesic), $\Omega(\partial_{u}) = \tfrac12\partial_{i}H\,e_{i}\wedge K$ and the projectors $\Lambda(x)$, $\Lambda(x')^{\dagger}$ of the parallel frame, $N_{0} = N_{1} = 0$ and $\mathrm{tr}(M^{-1}X) = 0$ exactly at three point pairs ($u > u'$, $u < u'$, central and off-central generators), while freezing $\Lambda(x)$, doubling $\partial_{u}\log\Delta$ or dropping the transport breaks it. (5) The same on the flat field with an expanding and shearing null hypersurface. (6) The rotor is the spin transport ($D_{T}R = 0$ along the geodesic tangent $T$), and $\Omega$ is checked against its coordinate reading (Christoffel symbols). In this gauge the connection enters the kernel only through the transport rotor (the term $K\,\Omega(\partial_{u})$ vanishes, $K(e_{i}\wedge K) = 0$), and the rotor is built from the same $\Omega(\partial_{u})$, so a change of its sign is detected; the projected transverse connection of Step 3(b) is tested on generic jets. (7) The mass and $V\log\sigma$ numerators vanish on the generator. (8) The model integrals of Steps 3 and 4.

\emph{Status.} (i) and (ii) are proved, given the imported leading Hadamard coefficient and $i0$ prescription and Theorem \claimref{synthesis.field-strength-cancels-from-the-leading-null-singularity-of-hadamard-dirac-states-on-null-hypersurfaces}(i), (ii); (iii) is proved from them. The restriction is the $\epsilon$-regularised one in which (HS) is posed; its agreement with the restriction defined through the characteristic-data embedding, used to pass from the parametrix to states, is not part of this result. Transport-absent.
